% !TeX program = xelatex
% !TeX options = -shell-escape -buf-size=500000
% !TeX encoding = UTF-8
%
% ╔══════════════════════════════════════════════════════════════════════════╗
% ║                                                                        ║
% ║   ANPG CONTEÚDO LOCAL — PROFESSIONAL DATA REPORT                       ║
% ║   Complete Data Export from Power BI Dashboard                         ║
% ║                                                                        ║
% ║   Source:     https://anpg.co.ao/conteudo-local/                       ║
% ║   Refresh:    2026-04-10                                               ║
% ║   Records:    81,215 rows · 3,595 companies · 20 categories            ║
% ║                                                                        ║
% ╚══════════════════════════════════════════════════════════════════════════╝
%
% ────────────────────────────────────────────────────────────────────────────
% DESIGN RATIONALE
% ────────────────────────────────────────────────────────────────────────────
%
% This template transforms a 1 300-page Power BI data export into a
% structured, navigable, professional report suitable for executive review,
% technical analysis, and archival.  Key design decisions:
%
% 1. DOCUMENT CLASS: report — not article.
%    article's \section is the top-level heading → inadequate for 1 300 pp.
%    report provides \part + \chapter + \section hierarchy.
%    openany: no blank even-sided pages (pure data document, not a novel).
%    twoside: alternating inner/outer margins for print / bindability;
%    running headers switch between recto and verso pages.
%
% 2. FOUR-PART STRUCTURE
%    The data naturally divides into four logical volumes:
%      Part I   — Overview & Legal Framework (chapters 1–4, ~160 lines)
%      Part II  — Service & Goods Taxonomy  (chapter 5, 20 sections)
%      Part III — Company Directory          (chapter 6, 1 massive table)
%      Part IV  — Company–Activity Cross-Ref (chapter 7, 803 sub-tables)
%    \part pages create visual "breathing room" in the 1 300-page document
%    and segment the TOC into navigable groups.
%
% 3. HEADING MAPPING (Markdown → LaTeX)
%    # (headingOne)   → suppressed    (title already on cover page)
%    ## (headingTwo)   → \chapter*    (7 main data sections)
%    ### (headingThree) → \section*   (categories, subsections)
%    #### (headingFour) → \subsection* (803 activity listings)
%    Star variants avoid double-numbering (markdown headings embed their
%    own numbers like "1. Key Statistics", "5. COMPLETE SERVICE…").
%    \addcontentsline ensures TOC and PDF bookmarks are still generated.
%
% 4. ASYMMETRIC MARGINS (twoside)
%    Inner margin (binding) = 2.8 cm; outer = 2.0 cm.
%    Professional print output if bound, while maximising column width
%    for data-dense tables.
%
% 5. TABLE STRATEGY
%    The markdown package renders pipe tables as longtable with booktabs.
%    Zebra striping via \rowcolors resets at each longtable (etoolbox hook).
%    Tables rendered at \small to maximise column width.
%    \sloppy + \emergencystretch handle long Portuguese company names.
%
% 6. COLOUR PALETTE — anchored to ANPG's institutional identity:
%    anpgBlue (#003D6B) — headings, structural elements, links
%    anpgGold (#D4A843) — decorative accents, part labels
%    tableStripe (#EBF0F5) — blue-tinted alternating table rows
%    Regime colours: green (E), amber (P), blue (C)
%
% 7. EXECUTIVE SUMMARY
%    A pre-data chapter with KPI boxes and pgfplots charts.
%    KPI values are \newcommand macros — easy to update manually or
%    by running  python preprocess-anpg-data.py  (generates kpi-values.tex).
%
% 8. PORTUGUESE LANGUAGE SUPPORT
%    babel[portuguese] ensures correct hyphenation (crucial for long
%    Portuguese words like "Decomissionamento") and date formatting.
%
% 9. TOC DEPTH = 1
%    Printed TOC shows \part + \chapter + \section (~56 entries).
%    Including \subsection would add 803+ entries → unusable.
%    PDF bookmarks show ALL levels (bookmarksdepth=4) for sidebar navigation.
%
% 10. FONTS
%    Times New Roman (body) — formal, government-standard, full Unicode.
%    Arial (headings) — clean contrast with body text.
%    Consolas (mono) — for NIF codes and technical identifiers.
%
% ALTERNATIVE WORKFLOWS (see also preprocess-anpg-data.py):
%    A. Keep \markdownInput (current, simplest).
%    B. Convert .md → .tex via Python for full per-table control
%       (landscape, sparklines, custom column widths).
%    C. Split .md into per-chapter files for selective compilation.
%
% ────────────────────────────────────────────────────────────────────────────
% BUILD INSTRUCTIONS
% ────────────────────────────────────────────────────────────────────────────
%
% Two passes required (TOC + bookmarks):
%
%   xelatex -buf-size=500000 --shell-escape anpg-local-content-full-data.tex
%   xelatex -buf-size=500000 --shell-escape anpg-local-content-full-data.tex
%
% NOTE: -buf-size=500000 is required because the markdown package
% generates the 3 596-row Company Directory table as a single
% 306K-character line, exceeding TeX's default 200K buffer.
%
% OPTIONAL — pre-extract KPIs (generates kpi-values.tex):
%   python preprocess-anpg-data.py
%
% DEPENDENCIES:
%   XeLaTeX · shell-escape · MiKTeX or TeX Live
%   System fonts: Times New Roman, Arial, Consolas
%   Source:       anpg-local-content-full-data.md (same directory)
%
% ────────────────────────────────────────────────────────────────────────────


% ══════════════════════════════════════════════════════════════════════════
%  DOCUMENT CLASS
% ══════════════════════════════════════════════════════════════════════════
\documentclass[
    a4paper,
    11pt,
    openany,        % no blank pages before chapters
    twoside,        % asymmetric margins for binding / recto-verso headers
]{report}


% ══════════════════════════════════════════════════════════════════════════
%  PAGE GEOMETRY  (asymmetric for binding)
% ══════════════════════════════════════════════════════════════════════════
\usepackage[
    a4paper,
    inner=2.8cm,        % wider binding margin
    outer=2.0cm,
    top=2.5cm,
    bottom=2.5cm,
    headheight=15pt,
    marginparwidth=1.5cm,
]{geometry}


% ══════════════════════════════════════════════════════════════════════════
%  LANGUAGE  (Portuguese primary — correct hyphenation & date formatting)
% ══════════════════════════════════════════════════════════════════════════
\usepackage[english,main=portuguese]{babel}


% ══════════════════════════════════════════════════════════════════════════
%  FONTS  (XeLaTeX / LuaLaTeX / pdfLaTeX fallback)
% ══════════════════════════════════════════════════════════════════════════
\usepackage{iftex}
\ifXeTeX
    \usepackage{fontspec}
    \defaultfontfeatures{Ligatures=TeX}
    \setmainfont{Times New Roman}
    \setsansfont{Arial}
    \setmonofont{Consolas}[Scale=0.88]
\else\ifLuaTeX
    \usepackage{fontspec}
    \defaultfontfeatures{Ligatures=TeX}
    \setmainfont{Times New Roman}
    \setsansfont{Arial}
    \setmonofont{Consolas}[Scale=0.88]
\else
    \usepackage[utf8]{inputenc}
    \usepackage[T1]{fontenc}
\fi\fi


% ══════════════════════════════════════════════════════════════════════════
%  COLOUR PALETTE
% ══════════════════════════════════════════════════════════════════════════
\usepackage[table,dvipsnames]{xcolor}

% ANPG corporate-inspired palette
\definecolor{anpgBlue}{HTML}{003D6B}        % deep petroleum blue
\definecolor{anpgDarkBlue}{HTML}{002244}     % part headings
\definecolor{anpgGold}{HTML}{D4A843}         % warm gold accent
\definecolor{anpgGray}{HTML}{F5F5F5}         % light background
\definecolor{tableStripe}{HTML}{EBF0F5}      % blue-tinted table stripe
\definecolor{linkBlue}{HTML}{1A5276}         % hyperlink tint
\definecolor{ruleGray}{HTML}{CCCCCC}         % decorative rules

% Regime indicator colours
\definecolor{regimeE}{HTML}{1B7340}          % Exclusividade — green
\definecolor{regimeP}{HTML}{B8860B}          % Preferência   — dark goldenrod
\definecolor{regimeC}{HTML}{2E6EA6}          % Concorrência  — steel blue

% KPI box colours
\definecolor{kpiBoxBg}{HTML}{F8FBFD}         % very light blue fill
\definecolor{kpiBoxBorder}{HTML}{B8D4E8}     % soft blue border


% ══════════════════════════════════════════════════════════════════════════
%  TABLES
% ══════════════════════════════════════════════════════════════════════════
\usepackage{longtable}
\usepackage{booktabs}
\usepackage{tabularx}
\usepackage{caption}
\usepackage{array}
\usepackage{multirow}
\usepackage{pdflscape}      % \begin{landscape} — rotated pages in PDF viewer

% Slightly open rows for data-dense pages
\renewcommand{\arraystretch}{1.20}

% Caption styling — small sans-serif, blue label
\captionsetup{
    font={small,sf},
    labelfont={bf,color=anpgBlue},
    skip=8pt,
    format=hang,
}

% Zebra striping — reset at the start of every longtable
\usepackage{etoolbox}
\AtBeginEnvironment{longtable}{%
    \rowcolors{2}{tableStripe}{white}%
    \small%
}


% ══════════════════════════════════════════════════════════════════════════
%  GRAPHICS & CHARTS
% ══════════════════════════════════════════════════════════════════════════
\usepackage{graphicx}
\usepackage{pgfplots}
\pgfplotsset{compat=1.18}
\usepackage{tikz}
\usetikzlibrary{calc,positioning}

% Sparklines note:  For mini in-cell charts, add \usepackage{sparklines}
% and pre-process the Markdown tables into LaTeX with the Python script
% (preprocess-anpg-data.py) to inject \begin{sparkline} environments.
% This is an optional refinement beyond the markdown-input workflow.


% ══════════════════════════════════════════════════════════════════════════
%  URLS & LINKS
% ══════════════════════════════════════════════════════════════════════════
\usepackage{xurl}          % breakable URLs at any character


% ══════════════════════════════════════════════════════════════════════════
%  MICROTYPOGRAPHY
% ══════════════════════════════════════════════════════════════════════════
\IfFileExists{microtype.sty}{%
    \usepackage[protrusion=true,expansion=false]{microtype}%
}{}


% ══════════════════════════════════════════════════════════════════════════
%  HEADERS & FOOTERS  (recto / verso differentiated)
% ══════════════════════════════════════════════════════════════════════════
\usepackage{fancyhdr}

% Body page style
\fancypagestyle{body}{%
    \fancyhf{}%
    \renewcommand{\headrulewidth}{0.4pt}%
    \renewcommand{\footrulewidth}{0pt}%
    % Even (verso/left): chapter title on left
    \fancyhead[LE]{\small\sffamily\textcolor{anpgBlue}{\leftmark}}%
    % Odd (recto/right): section title on right
    \fancyhead[RO]{\small\sffamily\textcolor{anpgBlue}{\rightmark}}%
    % Opposite corners: fixed label
    \fancyhead[RE,LO]{\small\sffamily\textcolor{anpgBlue}{ANPG Conteúdo Local}}%
    % Footer: centred page number + date stamp in corners
    \fancyfoot[C]{\small\sffamily\thepage}%
    \fancyfoot[LE,RO]{\footnotesize\sffamily\textcolor{ruleGray}{Abril 2026}}%
}

% Plain style — chapter opening pages (no header, centred page number)
\fancypagestyle{plain}{%
    \fancyhf{}%
    \renewcommand{\headrulewidth}{0pt}%
    \renewcommand{\footrulewidth}{0pt}%
    \fancyfoot[C]{\small\sffamily\thepage}%
}


% ══════════════════════════════════════════════════════════════════════════
%  SECTION HEADING STYLES
% ══════════════════════════════════════════════════════════════════════════
\usepackage{titlesec}

% Part — gold rules framing the part number, title in blue below
\titleformat{\part}[display]
    {\centering\normalfont\Huge\bfseries\sffamily\color{anpgBlue}}
    {\vspace*{3cm}{\color{anpgGold}\titlerule[2.5pt]}\vspace{8pt}%
     {\Large\textcolor{anpgGold}{Parte~\thepart}}\vspace{8pt}%
     {\color{anpgGold}\titlerule[2.5pt]}}
    {20pt}
    {}
    [\thispagestyle{empty}]

% Chapter — double rule + gold label
\titleformat{\chapter}[display]
    {\normalfont\huge\bfseries\sffamily\color{anpgBlue}}
    {\titlerule[2.5pt]\vspace{2pt}\titlerule[0.8pt]\vspace{6pt}%
     \Large\textcolor{anpgGold}{Capítulo~\thechapter}}
    {0pt}
    {\vspace{8pt}\Huge}
    [\vspace{6pt}{\color{ruleGray}\titlerule[0.8pt]}]

% Chapter* (unnumbered) — single rule, no "Capítulo N" label
\titleformat{name=\chapter,numberless}
    {\normalfont\huge\bfseries\sffamily\color{anpgBlue}}
    {}
    {0pt}
    {\vspace{8pt}}
    [\vspace{4pt}{\color{ruleGray}\titlerule[0.8pt]}]

% Section — blue with underline
\titleformat{\section}
    {\normalfont\Large\bfseries\sffamily\color{anpgBlue}}
    {\thesection}{0.8em}{}
    [{\color{ruleGray}\titlerule[0.4pt]}]

% Section* (unnumbered)
\titleformat{name=\section,numberless}
    {\normalfont\Large\bfseries\sffamily\color{anpgBlue}}
    {}
    {0pt}
    {}
    [{\color{ruleGray}\titlerule[0.4pt]}]

% Subsection — lighter blue
\titleformat{\subsection}
    {\normalfont\large\bfseries\sffamily\color{anpgBlue!80!black}}
    {\thesubsection}{0.7em}{}

% Subsection* (unnumbered)
\titleformat{name=\subsection,numberless}
    {\normalfont\large\bfseries\sffamily\color{anpgBlue!80!black}}
    {}
    {0pt}
    {}

% Subsubsection
\titleformat{\subsubsection}
    {\normalfont\normalsize\bfseries\sffamily}
    {\thesubsubsection}{0.6em}{}

% Spacing: {left}{before}{after}
\titlespacing*{\chapter}{0pt}{0pt}{24pt}
\titlespacing*{\section}{0pt}{20pt}{10pt}
\titlespacing*{\subsection}{0pt}{16pt}{8pt}
\titlespacing*{\subsubsection}{0pt}{12pt}{5pt}


% ══════════════════════════════════════════════════════════════════════════
%  TABLE OF CONTENTS STYLING
% ══════════════════════════════════════════════════════════════════════════
\usepackage{titletoc}

% Part — prominent, dark blue
\titlecontents{part}[0em]
    {\addvspace{18pt}\bfseries\Large\sffamily\color{anpgDarkBlue}}
    {\contentslabel{2.5em}}{}
    {}

% Chapter — bold blue with dot leaders
\titlecontents{chapter}[0em]
    {\addvspace{8pt}\bfseries\sffamily\color{anpgBlue}}
    {\contentslabel{1.8em}}{}
    {\titlerule*[8pt]{.}\contentspage}

% Section — indented, regular weight
\titlecontents{section}[1.5em]
    {\addvspace{2pt}\sffamily}
    {\contentslabel{2.8em}}{}
    {\titlerule*[6pt]{.}\contentspage}

% Subsection — further indented, small (not printed by default)
\titlecontents{subsection}[3.8em]
    {\small\sffamily}
    {\contentslabel{3.5em}}{}
    {\titlerule*[5pt]{.}\contentspage}

% Printed TOC depth: part + chapter + section (~56 entries)
\setcounter{tocdepth}{1}


% ══════════════════════════════════════════════════════════════════════════
%  LISTS
% ══════════════════════════════════════════════════════════════════════════
\usepackage{enumitem}
\setlist{nosep,leftmargin=1.5em}
\setlist[itemize]{label=\textcolor{anpgBlue}{\textbullet}}
\setlist[description]{style=nextline,leftmargin=2cm,font=\sffamily\bfseries\color{anpgBlue}}


% ══════════════════════════════════════════════════════════════════════════
%  HYPERLINKS & PDF METADATA
% ══════════════════════════════════════════════════════════════════════════
\usepackage[
    colorlinks=true,
    linkcolor=anpgBlue,
    urlcolor=linkBlue,
    citecolor=anpgBlue,
    breaklinks,
]{hyperref}
\usepackage{bookmark}

\hypersetup{
    pdftitle  = {ANPG Conteúdo Local — Relatório Completo de Dados},
    pdfauthor = {Extracção automatizada — Abril 2026},
    pdfsubject= {Registo, certificação e classificação de conteúdo local
                 de 3\,595 empresas no sector petrolífero angolano},
    pdfkeywords={ANPG, Conteúdo Local, Angola, petróleo, Power BI,
                 Decreto 86/18, Decreto 271/20},
}

% PDF bookmarks: show ALL levels for sidebar navigation
% (even though printed TOC only shows down to \section)
\bookmarksetup{depth=4}


% ══════════════════════════════════════════════════════════════════════════
%  MARKDOWN PACKAGE  (content import — needs --shell-escape)
% ══════════════════════════════════════════════════════════════════════════
\usepackage[pipeTables]{markdown}

% HEADING RENDERER CUSTOMISATION
%
% Strategy: use star (*) variants for all headings to avoid double-numbering
% (the Markdown source already embeds numbers: "## 1. Key Statistics", etc.).
% Manual \addcontentsline ensures TOC + PDF bookmarks are generated.
% \markboth / \markright feed the running headers (fancyhdr).
%
% \part boundaries are injected at strategic chapter indices to create
% the four-volume structure documented in the Design Rationale above.

\newcounter{mdchapter}

% # Title → suppressed (already on cover page / executive summary)
\def\markdownRendererHeadingOne#1{}

% ## → \chapter* with \part injection
\def\markdownRendererHeadingTwo#1{%
    \stepcounter{mdchapter}%
    \ifnum\value{mdchapter}=1\relax
        \part{Visão Geral e Enquadramento Legal}%
    \fi
    \ifnum\value{mdchapter}=5\relax
        \part{Taxonomia de Serviços e Bens}%
    \fi
    \ifnum\value{mdchapter}=6\relax
        \part{Directório de Empresas}%
    \fi
    \ifnum\value{mdchapter}=7\relax
        \part{Referência Cruzada Empresa–Actividade}%
    \fi
    \chapter*{#1}%
    \addcontentsline{toc}{chapter}{#1}%
    \markboth{#1}{#1}%
}

% ### → \section*
\def\markdownRendererHeadingThree#1{%
    \section*{#1}%
    \addcontentsline{toc}{section}{#1}%
    \markright{#1}%
}

% #### → \subsection* (not in printed TOC; appears in PDF bookmarks)
\def\markdownRendererHeadingFour#1{%
    \subsection*{#1}%
    \addcontentsline{toc}{subsection}{#1}%
}


% ══════════════════════════════════════════════════════════════════════════
%  KPI PLACEHOLDER COMMANDS
% ══════════════════════════════════════════════════════════════════════════
% Update these manually or regenerate with:
%   python preprocess-anpg-data.py   →   kpi-values.tex
% ──────────────────────────────────────────────────────────────────────────

% Registration pipeline
\newcommand{\kpiRegistered}{2\,534}
\newcommand{\kpiVisited}{1\,054}
\newcommand{\kpiCertified}{1\,199}
\newcommand{\kpiTotalDB}{3\,595}

% Data volume
\newcommand{\kpiMappings}{81\,215}
\newcommand{\kpiCategories}{20}
\newcommand{\kpiActivities}{397}

% Company types
\newcommand{\kpiSCA}{2\,076}
\newcommand{\kpiSE}{893}
\newcommand{\kpiSCDA}{624}

% Local content proxy
\newcommand{\kpiAngolan}{2\,662}
\newcommand{\kpiAngolanPct}{74{,}0\,\%}
\newcommand{\kpiForeignPct}{26{,}0\,\%}
\newcommand{\kpiCertifiedPct}{33{,}4\,\%}

% Top nationalities (for charts)
\newcommand{\kpiUK}{173}
\newcommand{\kpiUS}{100}
\newcommand{\kpiFR}{85}
\newcommand{\kpiUAE}{80}
\newcommand{\kpiPT}{69}

% Load auto-generated overrides if the Python script has been run
\IfFileExists{kpi-values.tex}{% ──────────────────────────────────────────────────────────────
% kpi-values.tex — auto-generated by preprocess-anpg-data.py
% Generated: 2026-04-11 01:29
% Re-run:  python preprocess-anpg-data.py
% ──────────────────────────────────────────────────────────────

\renewcommand{\kpiRegistered}{2\,534}
\renewcommand{\kpiVisited}{1\,054}
\renewcommand{\kpiCertified}{1\,199}
\renewcommand{\kpiTotalDB}{3\,595}
\renewcommand{\kpiMappings}{81\,215}
\renewcommand{\kpiSCA}{2\,076}
\renewcommand{\kpiSE}{893}
\renewcommand{\kpiSCDA}{624}
\renewcommand{\kpiAngolan}{2\,662}
\renewcommand{\kpiAngolanPct}{74.0\,\%}
\renewcommand{\kpiForeignPct}{26.0\,\%}
\renewcommand{\kpiCertifiedPct}{33.4\,\%}
\renewcommand{\kpiUK}{173}
\renewcommand{\kpiUS}{100}
\renewcommand{\kpiFR}{85}
\renewcommand{\kpiUAE}{80}
\renewcommand{\kpiPT}{69}

}{}


% ══════════════════════════════════════════════════════════════════════════
%  CUSTOM ENVIRONMENTS & COMMANDS
% ══════════════════════════════════════════════════════════════════════════

% KPI highlight box (used in Executive Summary dashboard grid)
\newcommand{\kpibox}[3]{%
    % #1 = number, #2 = label, #3 = sublabel
    \begin{tikzpicture}
    \node[
        draw=kpiBoxBorder,
        fill=kpiBoxBg,
        rounded corners=4pt,
        minimum width=4.0cm,
        minimum height=2.4cm,
        inner sep=8pt,
        align=center,
    ] {
        {\fontsize{22}{26}\selectfont\bfseries\sffamily\textcolor{anpgBlue}{#1}}\\[4pt]
        {\small\sffamily #2}\\[1pt]
        {\scriptsize\sffamily\textcolor{gray}{#3}}
    };
    \end{tikzpicture}%
}

% Regime legend (inline colour-coded labels)
\newcommand{\regE}{\textcolor{regimeE}{\textbf{[E]}}}
\newcommand{\regP}{\textcolor{regimeP}{\textbf{[P]}}}
\newcommand{\regC}{\textcolor{regimeC}{\textbf{[C]}}}


% ══════════════════════════════════════════════════════════════════════════
%  PARAGRAPH & LINE-BREAKING TUNING
% ══════════════════════════════════════════════════════════════════════════
\sloppy
\setlength{\parindent}{0pt}
\setlength{\parskip}{4pt plus 1pt minus 1pt}
\setlength{\emergencystretch}{3em}


% ╔══════════════════════════════════════════════════════════════════════════╗
% ║                         DOCUMENT BEGINS                                ║
% ╚══════════════════════════════════════════════════════════════════════════╝

\begin{document}


% ══════════════════════════════════════════════════════════════════════════
%  COVER PAGE
% ══════════════════════════════════════════════════════════════════════════
\begin{titlepage}
\newgeometry{margin=2.5cm}
\thispagestyle{empty}

\begin{center}

\vspace*{1.5cm}

% ── Top decorative rules ──
{\color{anpgBlue}\rule{\linewidth}{3pt}}
\vspace{2pt}
{\color{anpgGold}\rule{\linewidth}{1pt}}

\vspace{1.8cm}

% ── LOGO PLACEHOLDER ──
% Replace 'anpg-logo.png' with the actual ANPG logo file.
% Uncomment when the logo is available:
% \includegraphics[width=5cm]{anpg-logo.png}\par\vspace{0.8cm}

{\fontsize{42}{48}\selectfont\bfseries\sffamily\textcolor{anpgBlue}{ANPG}}

\vspace{0.3cm}

{\fontsize{24}{30}\selectfont\sffamily\textcolor{anpgBlue}{Conteúdo Local}}

\vspace{0.8cm}

{\color{anpgGold}\rule{0.5\linewidth}{2pt}}

\vspace{1cm}

{\Large\sffamily Relatório Completo de Dados\\[6pt]
 Extraído do Painel Power~BI}

\vspace{2cm}

% ── Key figures box ──
\fcolorbox{anpgBlue}{kpiBoxBg}{%
\begin{minipage}{0.80\linewidth}
\centering\sffamily
\vspace{14pt}
{\large\bfseries\textcolor{anpgBlue}{Dados em Destaque}}\\[12pt]
\begin{tabular}{r@{\hspace{10pt}}l@{\qquad}r@{\hspace{10pt}}l}
    \textcolor{anpgBlue}{\textbf{\kpiTotalDB}}   & empresas no sistema &
    \textcolor{anpgBlue}{\textbf{\kpiCategories}} & categorias de serviços \\[4pt]
    \textcolor{anpgBlue}{\textbf{\kpiMappings}}   & mapeamentos empresa–actividade &
    \textcolor{anpgBlue}{\textbf{\kpiActivities}} & actividades distintas \\[4pt]
    \textcolor{anpgBlue}{\textbf{\kpiRegistered}} & empresas registadas &
    \textcolor{anpgBlue}{\textbf{\kpiCertified}}  & empresas certificadas \\
\end{tabular}
\vspace{14pt}
\end{minipage}%
}

\vspace{1.8cm}

{\large\sffamily Extracção automatizada — Abril 2026}

\vspace{0.4cm}

{\sffamily\textit{Última actualização dos dados: 10 de Abril de 2026}}

\vspace{0.6cm}

{\small\sffamily Fonte: \url{https://anpg.co.ao/conteudo-local/}}

\vfill

% ── Bottom decorative rules ──
{\color{anpgGold}\rule{\linewidth}{1pt}}
\vspace{2pt}
{\color{anpgBlue}\rule{\linewidth}{3pt}}

\vspace{0.6cm}

{\footnotesize\sffamily\textcolor{gray}{%
    Base legal: Decreto Presidencial 86/18 (Conteúdo Local)
    \quad·\quad Decreto Presidencial Nº\,271/20 (Registo e Certificação)}}

\end{center}

\restoregeometry
\end{titlepage}


% ══════════════════════════════════════════════════════════════════════════
%  FRONT MATTER  (roman page numbers — pre-data content)
% ══════════════════════════════════════════════════════════════════════════
\pagenumbering{roman}
\pagestyle{body}


% ── Legal Notice ─────────────────────────────────────────────────────────
\chapter*{Nota Legal}
\addcontentsline{toc}{chapter}{Nota Legal}
\markboth{Nota Legal}{Nota Legal}

Este documento foi gerado automaticamente a partir dos dados públicos
disponíveis no painel Power~BI de Conteúdo Local da ANPG
(\url{https://anpg.co.ao/conteudo-local/}).

Os dados aqui apresentados são de natureza pública e destinam-se a fins
de análise, investigação e referência técnica.
A informação é apresentada tal como foi extraída, sem modificações ao
conteúdo original.

\begin{description}
    \item[Decreto Presidencial 86/18]
        Regime Jurídico do Conteúdo Local no Sector Petrolífero.
    \item[Decreto Presidencial Nº\,271/20]
        Registo e Certificação de Empresas para o Sector Petrolífero em Angola.
\end{description}

A ANPG publica anualmente as listas de bens e serviços a contratar pelos
operadores e associados sob os regimes de:

\begin{itemize}
    \item \regE{} \textbf{Exclusividade} — contratação exclusiva a empresas
          locais/angolanas
    \item \regP{} \textbf{Preferência} — preferência para empresas locais
    \item \regC{} \textbf{Concorrência} — concurso aberto
\end{itemize}

\clearpage


% ── How to Read This Document ────────────────────────────────────────────
\chapter*{Como Utilizar Este Documento}
\addcontentsline{toc}{chapter}{Como Utilizar Este Documento}
\markboth{Como Utilizar Este Documento}{Como Utilizar Este Documento}

Este relatório está organizado em \textbf{quatro partes}:

\begin{description}
    \item[Parte~I — Visão Geral e Enquadramento Legal]
        Estatísticas-chave, enquadramento legal, estrutura do formulário de
        registo e modelo de dados do Power~BI.  Ideal para compreender o
        contexto antes de consultar os dados detalhados.

    \item[Parte~II — Taxonomia de Serviços e Bens]
        As \kpiCategories{} categorias de serviços e bens com as respectivas
        actividades, tipo (serviço/bem) e regime (Exclusividade, Preferência
        ou Concorrência).  Utilize esta parte para identificar nichos de
        mercado e oportunidades de conteúdo local.

    \item[Parte~III — Directório de Empresas]
        Listagem completa das \kpiTotalDB{} empresas, ordenadas
        alfabeticamente, com NIF, tipo de sociedade, estado de certificação
        e nacionalidade.

    \item[Parte~IV — Referência Cruzada Empresa–Actividade]
        Os \kpiMappings{} registos que ligam cada empresa às actividades
        em que está inscrita.  Organizado por Categoria → Actividade →
        Empresas.  Utilize esta parte para pesquisar concorrentes ou
        parceiros por área de serviço.
\end{description}

\medskip

\noindent\textbf{Navegação:}
Utilize o \textbf{Índice} (página seguinte) para navegar entre partes e
capítulos.  No leitor de PDF, os \textbf{marcadores laterais} (bookmarks)
incluem todos os níveis, permitindo saltar directamente para qualquer
actividade individual.

\clearpage


% ── Executive Summary ────────────────────────────────────────────────────
\chapter*{Sumário Executivo}
\addcontentsline{toc}{chapter}{Sumário Executivo}
\markboth{Sumário Executivo}{Sumário Executivo}

% ─── KPI Dashboard: 6 key metrics in a 3×2 grid ───
\begin{center}
\vspace{0.3cm}

\kpibox{\kpiTotalDB}{Empresas no Sistema}{total na base de dados}%
\hfill
\kpibox{\kpiRegistered}{Registadas}{\kpiAngolanPct{} angolanas}%
\hfill
\kpibox{\kpiCertified}{Certificadas}{\kpiCertifiedPct{} do total}%

\vspace{0.8cm}

\kpibox{\kpiMappings}{Mapeamentos}{empresa–actividade}%
\hfill
\kpibox{\kpiCategories}{Categorias}{de serviços e bens}%
\hfill
\kpibox{\kpiActivities}{Actividades}{distintas registadas}%

\vspace{0.3cm}
\end{center}


% ─── Chart 1: Company Status Breakdown ───
\vspace{0.8cm}
\noindent\textbf{\sffamily\large\textcolor{anpgBlue}{%
    Distribuição por Estado de Certificação}}

\vspace{0.5cm}

\begin{center}
\begin{tikzpicture}
\begin{axis}[
    xbar,
    width=0.80\textwidth,
    height=7cm,
    bar width=12pt,
    xlabel={\sffamily Número de empresas},
    xlabel style={font=\small\sffamily},
    xmin=0, xmax=1850,
    symbolic y coords={%
        {Outros},
        {Recusado},
        {Condicionado},
        {Aprovado},
        {Aprovado (Int.)},
        {Certificado}%
    },
    ytick=data,
    yticklabel style={font=\small\sffamily},
    xticklabel style={font=\small\sffamily},
    nodes near coords,
    nodes near coords style={font=\small\sffamily\bfseries,color=anpgBlue},
    every axis plot/.append style={fill=anpgBlue!70,draw=anpgBlue},
    enlarge y limits=0.15,
    grid=major,
    grid style={gray!20},
    axis line style={gray!50},
]
\addplot coordinates {
    (130,{Outros})
    (114,{Recusado})
    (430,{Condicionado})
    (475,{Aprovado})
    (859,{Aprovado (Int.)})
    (1587,{Certificado})
};
\end{axis}
\end{tikzpicture}
\end{center}

{\small\sffamily\textcolor{gray}{%
    «Outros» agrupa: Pré~Triagem (112), Triagem (11),
    Em~Certificação (4+2+1).}}


% ─── Chart 2: Company Types ───
\vspace{0.8cm}
\noindent\textbf{\sffamily\large\textcolor{anpgBlue}{%
    Distribuição por Tipo de Sociedade}}

\vspace{0.5cm}

\begin{center}
\begin{tikzpicture}
\begin{axis}[
    xbar,
    width=0.62\textwidth,
    height=4cm,
    bar width=16pt,
    xlabel={\sffamily Número de empresas},
    xlabel style={font=\small\sffamily},
    xmin=0,
    symbolic y coords={SCDA, SE, SCA},
    ytick=data,
    yticklabel style={font=\sffamily},
    xticklabel style={font=\small\sffamily},
    nodes near coords,
    nodes near coords style={font=\sffamily\bfseries,color=anpgBlue},
    every axis plot/.append style={fill=anpgGold!80!anpgBlue,draw=anpgBlue!60},
    enlarge y limits=0.35,
    axis line style={gray!50},
]
\addplot coordinates {
    (624,SCDA)
    (893,SE)
    (2076,SCA)
};
\end{axis}
\end{tikzpicture}
\end{center}

\begin{center}
\small\sffamily
\textbf{SCA} — Sociedade Comercial Angolana \quad
\textbf{SE} — Sociedade Estrangeira \quad
\textbf{SCDA} — Sociedade Comercial de Direito Angolano
\end{center}


% ─── Chart 3: Top 10 Service Categories by Volume ───
\vspace{0.8cm}
\noindent\textbf{\sffamily\large\textcolor{anpgBlue}{%
    Top~10 Categorias por Volume de Registos}}

\vspace{0.5cm}

\begin{center}
\begin{tikzpicture}
\begin{axis}[
    xbar,
    width=0.80\textwidth,
    height=9.5cm,
    bar width=10pt,
    xlabel={\sffamily Mapeamentos empresa–actividade},
    xlabel style={font=\small\sffamily},
    xmin=0,
    symbolic y coords={%
        {Formação},
        {Suporte Operações},
        {Engenharia},
        {Inspecção/Testes},
        {Consultoria},
        {Fabricação},
        {Logísticos},
        {HSE},
        {Perfuração},
        {Apoio Geral},
        {TIC}%
    },
    ytick=data,
    yticklabel style={font=\scriptsize\sffamily},
    xticklabel style={font=\small\sffamily},
    nodes near coords,
    nodes near coords style={font=\tiny\sffamily\bfseries,color=anpgBlue},
    every axis plot/.append style={fill=anpgBlue!55,draw=anpgBlue!80},
    enlarge y limits=0.08,
    grid=major,
    grid style={gray!15},
    axis line style={gray!50},
]
\addplot coordinates {
    (4788,{Formação})
    (4490,{Suporte Operações})
    (5492,{Engenharia})
    (5943,{Inspecção/Testes})
    (3987,{Consultoria})
    (3985,{Fabricação})
    (3882,{Logísticos})
    (3892,{HSE})
    (8141,{Perfuração})
    (10433,{Apoio Geral})
    (10441,{TIC})
};
\end{axis}
\end{tikzpicture}
\end{center}

{\small\sffamily\textcolor{gray}{%
    As 20 categorias completas com totais detalhados encontram-se na
    Parte~II (Capítulo~5).}}


% ─── Local Content Overview ───
\vspace{1cm}
\noindent\textbf{\sffamily\large\textcolor{anpgBlue}{%
    Indicadores de Conteúdo Local}}

\vspace{0.4cm}

\begin{center}
\begin{tabular}{lrl}
\toprule
\textbf{Indicador} & \textbf{Valor} & \textbf{Nota} \\
\midrule
Empresas angolanas (SCA+SCDA)
    & \kpiAngolan{}
    & \kpiAngolanPct{} do total \\
Empresas estrangeiras (SE)
    & \kpiSE{}
    & \kpiForeignPct{} do total \\
Empresas certificadas
    & \kpiCertified{}
    & \kpiCertifiedPct{} taxa de certificação \\
Categorias de serviços
    & \kpiCategories{}
    & abrangendo \kpiActivities{} actividades \\
\bottomrule
\end{tabular}
\end{center}

\vspace{0.5cm}

{\small\sffamily\textcolor{gray}{%
    Nota: os indicadores acima são derivados da contagem directa dos
    registos no Power~BI.  As métricas de receita, colaboradores e
    percentagem de conteúdo local por empresa estão disponíveis nas
    tabelas «Factual de Conteúdo Local» e «Factual Contabilidade» do
    modelo de dados (Capítulo~4).}}

\clearpage


% ── Table of Contents ────────────────────────────────────────────────────
\tableofcontents

\clearpage


% ══════════════════════════════════════════════════════════════════════════
%  MAIN CONTENT  (arabic page numbers — raw data from Power BI)
% ══════════════════════════════════════════════════════════════════════════
\pagenumbering{arabic}
\setcounter{page}{1}

% ── DATA INPUT ───────────────────────────────────────────────────────────
% The Markdown file is ingested below.  Custom heading renderers (defined
% in the preamble) map ## → \chapter*, ### → \section*, #### → \subsection*
% and inject \part boundaries at chapters 1, 5, 6, and 7.
%
% ALTERNATIVE WORKFLOW:
% If you convert the Markdown to LaTeX fragments using the preprocessing
% script, replace the line below with:
%   \input{anpg-local-content-data.tex}
% and gain per-table control (landscape, sparklines, custom column widths).
% ─────────────────────────────────────────────────────────────────────────

\markdownInput{anpg-local-content-full-data.md}


\end{document}
